\documentclass{article}
\usepackage[english]{babel}
\usepackage[toc,page]{appendix}
\usepackage[letterpaper,top=2cm,bottom=2cm,left=3cm,right=3cm,marginparwidth=1.75cm]{geometry}
\usepackage{enumitem}
\usepackage{amsmath}
\usepackage[colorlinks=true, allcolors=blue]{hyperref}
\usepackage{minted}
\usepackage{float}
\usepackage{graphicx}
\usepackage{multirow}
\usepackage{subcaption}
\usepackage{booktabs}
\usepackage[utf8]{inputenc}
\usepackage{textgreek}
\usepackage{booktabs}


\title{Making a CPM of Collective Migration\\[0.3em] \large Report on Assignment 1C}
\author{Aaron Bussche \and  David Magaram\\[0.3em] Group 28 \and Jason Kasdiran}
\date{\today}

\begin{document}
\maketitle




\section{Introduction}

Collective cell migration concerns the movement of a group of cells in response to environmental or chemical cues, and is an emergent property of many interacting individual cells. In this phenomenon, cells in the back of a group can push the front cells to move quicker than they otherwise would have. Cells in the back can place pixels inside cells that are in front of them, thereby putting the front cells 'under volume'. This results in the front cells placing new pixels in front of them, resulting in faster cell movement. We can explore how individual cells interacting with their environment give rise to such collective behavior, by using the Cellular Potts Model (CPM).

Obstacles represent some environmental constraints that affect collective cell migration by guiding the movement of cells or disrupting it. By varying the number of obstacles, changes in the collective cell migration can be observed. Investigating these effects helps reveal the underlying mechanisms that guide cell and group behaviors. This experiment aims to explore how obstacles change the collective migration behavior of cells by observing simulations with 0, 9, 25, and 49 obstacles.

The goal of this project is to use CPM to investigate the impact of obstacles on collective cell migration. 



\section{Methodology}
The Artistoo framework was used to implement the two-dimensional CPM simulations. The model was defined on a square grid of 200x200 pixels with wrapped dimensions. In addition to the background, there were two types of cells in the simulations: moving cells and obstacles. Obstacles were implemented by creating a new type of cell with different parameters in order to distinguish the two types of cells and to give them desired properties. All parameters used in this project are displayed in Table 1. 

Obstacles were initialized with a target volume ($V$) of 100 pixels and volume constraint strength ($\lambda_V$) of 500 to ensure they are static and do not deform. Their perimeter ($P$) was set to 50 with a strong perimeter constraint strength ($\lambda_P$) of 200, keeping their shape strongly constrained. To further prevent any movement, their activity ($\lambda_{\text{act}}, \text{Max}_{\text{act}}$) was set to 0. Additionally, the adhesion between obstacles and moving cells ($J_{\text{obstacle-cell}}$) was set to 1000 to ensure moving cells do not stick to obstacles or penetrate them. Lastly, obstacles were initialized with obstacle padding. This padding is the number of pixels between the border and obstacles. All obstacles within this bounding box are equally spaced in a grid-like pattern. 


On the other hand, moving cells had completely different parameters that would allow them to move and explore their environment, reproducing collective cell migration. These parameters were directly taken from the self-study exercise 1.3. In total, there were 120 moving cells on the grid. 120 was chosen because it does not clutter the screen and it allows for the observation of collective cell migration. Each moving cell was initialized with V of 200 pixels, making them twice the size of the obstacles, and $\lambda_V$ of 50 to allow them moderate compressibility. Their P was set to 180 and $\lambda_P$ to 2 to allow them to take on irregular shapes during their migration. Their movement, and the active migration, was controlled using activity constraints with $\lambda_{act}$ of 200 and $\text{Max}_{\text{act}}$ 20. These parameters allowed the reproduction of collective migration. 

\begin{table}[h!]
\centering
\begin{tabular}{lccc}
\hline
\textbf{Parameter} & \textbf{Moving Cells} & \textbf{Obstacles} \\
\hline
Temperature (T) & 20 & 20 \\
Adhesion cell–matrix ($J_{\text{obstacle-matrix}}$) & 20 & 20 \\
Adhesion cell–cell ($J_{\text{cell-cell}}$) & 0 & 0 \\
Adhesion obstacle–cell ($J_{\text{obstacle-cell}}$) & 1000 & 1000 \\
Target volume (V) & 200 & 100 \\
Volume constraint ($\lambda_V$) & 50 & 500 \\
Target perimeter (P) & 180 & 50 \\
Perimeter constraint ($\lambda_P$) & 2 & 200 \\
Activity strength ($\lambda_{\text{act}}$)& 200 & 0 \\
Activity memory ($\text{Max}_{\text{act}}$)  & 20 & 0 \\
Number of cells & 120 & 0, 9, 25, 49 \\
Obstacle padding  & -- & 10 \\
\hline
\end{tabular}
\caption{CPM parameters for moving cells and obstacles.}
\end{table}

There were 4 experiments in total. All were performed under the same conditions and with identical parameters, with the exception of the number of obstacles. In order to investigate how obstacles change the collective migration behavior of cells, 4 obstacle configurations were explored: 0, 9, 25, and 49. Obstacles placed on the lattice were equally spaced from each other, forming a grid-like structure. Additionally, throughout the experiments all simulations had temperature (T) 20, which controls the stochasticity of the system, allowing cells to deform and migrate. Each simulation was initialized by placing the obstacles first, and then randomly seeding moving cells. Each simulation ran for 600 Monte Carlo steps, which took about 10 minutes in wall-clock time at 1 frame per second, and the observations and results are reported in the Results section. 

To get quantitative data, a Python script was used to calculate the average speed and distance traversed, as well as find the fastest and slowest cells, along with the standard deviation in speeds. It also calculates the Vicsek order parameter for each timestep and takes the mean, giving us data on how stable the alignment of moving cells is. This allows us to analyze the cell movements for both environments with and without obstacles in a quantitative, objective way, and base our qualitative analysis on facts. 



\section{Results}


\section{Conclusion}
This study demonstrated that ...


\section{Discussion}
This experiment explores how different obstacle configurations impact collective cell migration. While it achieves its goal, there are several limitations that were not addressed. 

All obstacle configurations were equally spaced and placed in a grid-like pattern. Because of this, the environment did not change much between obstacle configurations. Moreover, only 4 obstacle configurations were explored, all with the same layout. 
Furthermore, most parameter values were taken from the self-study exercise 1.3 and no other parameter configurations were explored. 
It is recommended that future studies explore different parameter configurations, irregular obstacle placement, and dynamic obstacles. 

\newpage

\begin{appendices}

\end{appendices}

\end{document}